\documentclass[11pt]{article}
\usepackage{amsmath,amssymb}
\usepackage{geometry}
\geometry{margin=1in}

\title{Hopkins-Style Supernova Coupling on a Cartesian Mesh}
\author{}
\date{}

\begin{document}
\maketitle

\section{Overview}
This note describes a Cartesian-mesh implementation of the Hopkins et al. mechanical
supernova (SN) coupling, adapted to a fixed-grid particle--mesh workflow with
order-independent multi-SN updates. The method injects ejecta mass, momentum, and
energy into mesh cells using vector weights that enforce isotropy and conservation,
then sets the coupled momentum $p_{0}$ from energy conservation and a terminal
momentum cap.

\section{Notation}
We use cell-centered conserved variables for each mesh cell $b$ with volume $V_{b}$:
mass $m_{b}=\rho_{b}V_{b}$, momentum $\mathbf p_{b}=\rho_{b}\mathbf v_{b}V_{b}$, total
energy $E_{b}$, and internal energy $U_{b}$. For a supernova event (star particle) $a$,
let $\mathbf x_{a}$ and $\mathbf v_{a}$ be position and velocity, with ejecta mass
$m_{\rm ej}$ and total ejecta energy $E_{\rm ej}$ defined in the star frame. We also
use the relative velocity $\mathbf v_{ba}=\mathbf v_{b}-\mathbf v_{a}$.

\section{Kernel and Vector Weights on a Cartesian Mesh}
Choose a kernel radius $r_{K}$ and include all cells $b$ with
$|\mathbf x_{b}-\mathbf x_{a}| \le r_{K}$. Define the cell bounds
$(x_{\min},x_{\max})$, $(y_{\min},y_{\max})$, $(z_{\min},z_{\max})$
relative to $\mathbf x_{a}$.

We compute the exact solid angle of the (axis-aligned) rectangular cell as the
sum of the solid angles of its visible faces (at most three). For a rectangle
lying in the plane $z=\text{const}$ with corners $(x_{1},y_{1})$ and $(x_{2},y_{2})$
as seen from the origin, the solid angle is
\begin{equation}
\Omega_{\rm rect}(x_{1},x_{2},y_{1},y_{2},z) =
\tan^{-1}\!\left(\frac{x_{1}y_{1}}{z\,r_{11}}\right)
- \tan^{-1}\!\left(\frac{x_{1}y_{2}}{z\,r_{12}}\right)
- \tan^{-1}\!\left(\frac{x_{2}y_{1}}{z\,r_{21}}\right)
+ \tan^{-1}\!\left(\frac{x_{2}y_{2}}{z\,r_{22}}\right),
\end{equation}
where $r_{ij}=\sqrt{x_{i}^{2}+y_{j}^{2}+z^{2}}$ and $z$ is the (signed) distance
to the rectangle plane. We use $|z|$ in the formula to keep $\Omega_{\rm rect}>0$.

For each cell, include the visible faces:
\begin{equation}
\Omega_{b} =
\Omega_{x} + \Omega_{y} + \Omega_{z},
\end{equation}
with
\begin{align}
\Omega_{x} &=
\begin{cases}
\Omega_{\rm rect}(y_{\min},y_{\max},z_{\min},z_{\max},x_{\min}), & x_{\min} > 0, \\
\Omega_{\rm rect}(y_{\min},y_{\max},z_{\min},z_{\max},x_{\max}), & x_{\max} < 0, \\
0, & \text{otherwise,}
\end{cases} \\
\Omega_{y} &=
\begin{cases}
\Omega_{\rm rect}(x_{\min},x_{\max},z_{\min},z_{\max},y_{\min}), & y_{\min} > 0, \\
\Omega_{\rm rect}(x_{\min},x_{\max},z_{\min},z_{\max},y_{\max}), & y_{\max} < 0, \\
0, & \text{otherwise,}
\end{cases} \\
\Omega_{z} &=
\begin{cases}
\Omega_{\rm rect}(x_{\min},x_{\max},y_{\min},y_{\max},z_{\min}), & z_{\min} > 0, \\
\Omega_{\rm rect}(x_{\min},x_{\max},y_{\min},y_{\max},z_{\max}), & z_{\max} < 0, \\
0, & \text{otherwise.}
\end{cases}
\end{align}
The solid-angle weight is then
\begin{equation}
\omega_{b} = \frac{\Omega_{b}}{4\pi}.
\end{equation}
Define a raw vector weight
\begin{equation}
\mathbf q_{b} = \omega_{b}\,\hat{\mathbf r}.
\end{equation}
To enforce exact linear-momentum conservation (zero net vector weight), remove the
dipole component:
\begin{equation}
\mathbf Q = \sum_{b}\mathbf q_{b}, \qquad
\mathbf q^{\prime}_{b} = \mathbf q_{b} - \frac{\omega_{b}}{\sum_{c}\omega_{c}}\,\mathbf Q.
\end{equation}
Normalize to obtain the Hopkins-style vector weights:
\begin{equation}
\bar{\mathbf w}_{b} = \frac{\mathbf q^{\prime}_{b}}{\sum_{c}|\mathbf q^{\prime}_{c}|}, \qquad
\sum_{b}\bar{\mathbf w}_{b}=\mathbf 0,\quad \sum_{b}|\bar{\mathbf w}_{b}|=1.
\end{equation}

For a uniform Cartesian mesh with $\Delta x=\Delta y=\Delta z$, a source at the
center of a cell, and a one-cell-thick kernel (the 26 neighbors in a $3^{3}$
stencil), the unnormalized solid-angle weights $\omega_{b}=\Omega_{b}/(4\pi)$
take only three distinct values, depending on the neighbor class. The table
below lists the raw $\omega_{b}$ and the normalized $\bar\omega_{b}$
obtained by dividing by the sum over the 26 cells. These values are specific to
this kernel choice; larger kernels require re-normalization.

\begin{center}
\begin{tabular}{lccc}
\hline
Offset class $(|i|,|j|,|k|)$ & Count & $\omega_{b}$ & $\bar\omega_{b}$ (3x3x3) \\
\hline
$(0,0,1)$ (face neighbor) & 6  & 0.1666666667 & 0.0779325222 \\
$(0,1,1)$ (edge neighbor) & 12 & 0.0673911931 & 0.0315117939 \\
$(1,1,1)$ (corner neighbor) & 8  & 0.0412384905 & 0.0192829174 \\
\hline
\end{tabular}
\end{center}

Define the mass-loading factor
\begin{equation}
w^{\prime}_{b} \equiv \frac{|\bar{\mathbf w}_{b}|}{1 + \Delta m_{b}/m_{b}},
\qquad
\Delta m_{b} \equiv |\bar{\mathbf w}_{b}|\,m_{\rm ej}.
\end{equation}

\section{Single-Event Coupling}
\subsection{Velocity and mass scalars}
Compute the effective radial velocity and coupled mass (Hopkins et al. notation):
\begin{equation}
\langle v_{r}\rangle = \sum_{b} w^{\prime}_{b}\,\mathbf v_{ba}\cdot\frac{\bar{\mathbf w}_{b}}{|\bar{\mathbf w}_{b}|},
\end{equation}
\begin{equation}
\langle M_{\rm coupled}\rangle = \left[\sum_{b}\frac{w^{\prime}_{b}|\bar{\mathbf w}_{b}|}{m_{b}}\right]^{-1}.
\end{equation}
Define
\begin{equation}
\beta_{1} = \frac{\langle v_{r}\rangle}{v_{\rm ej}}, \qquad
\beta_{2} = \frac{m_{\rm ej}}{\langle M_{\rm coupled}\rangle},
\end{equation}
where $v_{\rm ej}$ is the ejecta speed implied by the available kinetic energy.

\subsection{Energy partition and coupled momentum}
Let the total available energy (in the star frame) be
\begin{equation}
\mathcal E_{a} = E_{\rm ej} + \frac{1}{2}m_{\rm ej}\sum_{b} w^{\prime}_{b} |\mathbf v_{ba}|^{2}.
\end{equation}
Choose a fixed kinetic fraction $f_{\rm kin}^{0}$ (e.g. $0.28$), then
\begin{equation}
\epsilon = f_{\rm kin}^{0}\,\mathcal E_{a}.
\end{equation}
Define
\begin{equation}
\psi = \frac{\sqrt{\beta_{2} + \beta_{1}^{2}} - \beta_{1}}{\beta_{2}}.
\end{equation}
Given a terminal momentum $p_{\rm term}$, set
\begin{equation}
\chi = \min\left[1,\ \frac{p_{\rm term}}{\psi \sqrt{2\epsilon m_{\rm ej}}}\right], \qquad
p_{0} = \psi\,\chi\,\sqrt{2\epsilon m_{\rm ej}}.
\end{equation}
The non-kinetic energy fraction is then
\begin{equation}
\frac{U_{a}}{\mathcal E_{a}} =
1 - \left[(\psi\chi)^{2}\beta_{2} + 2(\psi\chi)\beta_{1}\right]\frac{\epsilon}{\mathcal E_{a}}.
\end{equation}

\subsection{Cell updates for a single event}
Define the momentum increment (in the lab frame)
\begin{equation}
\Delta \mathbf p_{b} = \Delta m_{b}\,\mathbf v_{a} + \bar{\mathbf w}_{b}\,p_{0}.
\end{equation}
Deposit internal energy and update total energy consistently:
\begin{equation}
\Delta U_{b} = |\bar{\mathbf w}_{b}|\,U_{a},
\end{equation}
\begin{equation}
\Delta E_{b} = \Delta U_{b}
 + \frac{|\mathbf p_{b}+\Delta \mathbf p_{b}|^{2}}{2(m_{b}+\Delta m_{b})}
 - \frac{|\mathbf p_{b}|^{2}}{2m_{b}}.
\end{equation}
Update the conserved variables in each cell:
\begin{equation}
m_{b} \leftarrow m_{b} + \Delta m_{b},\quad
\mathbf p_{b} \leftarrow \mathbf p_{b} + \Delta \mathbf p_{b},\quad
E_{b} \leftarrow E_{b} + \Delta E_{b},\quad
U_{b} \leftarrow U_{b} + \Delta U_{b}.
\end{equation}
When multiple events are applied in a single timestep, total energy should not
be accumulated by summing $\Delta E_{b}$, because the kinetic-energy contribution
depends on the \emph{summed} momentum and can cancel. A consistent approach is to
sum $\Delta \rho$, $\Delta \mathbf p$, and $\Delta U$ (or other non-kinetic
reservoirs), then reconstruct total energy in the simulation frame from
$E=U+|\mathbf p|^{2}/(2\rho)$ after the sum.

\section{Order-Independent Multi-SN Algorithm}
To preserve the GPU-friendly, order-independent behavior of the particle-mesh
method, compute each SN event using the pre-feedback state and accumulate
changes into a buffer mesh:
\begin{enumerate}
\item For each event $a$, compute $\bar{\mathbf w}_{b}$, $\psi$, $\chi$, $p_{0}$, and
the per-cell increments $(\Delta m_{b}, \Delta \mathbf p_{b}, \Delta U_{b})$
using the state at time $t^{n}$ only.
\item Accumulate all event contributions into buffer fields
$\Delta m^{\rm buf}$, $\Delta \mathbf p^{\rm buf}$, $\Delta U^{\rm buf}$
using atomic adds.
\item Communicate and sum buffer ghost zones across MPI ranks.
\item (Optional) apply a cell-wise positivity limiter to the summed increments
if needed for robustness; otherwise skip, since the Hopkins coupling is
energy-conserving per event.
\item Apply the summed increments to $\rho$, $\mathbf p$, and $U$.
\item Reconstruct total energy in the simulation frame from
$E=U+|\mathbf p|^{2}/(2\rho)$ after the sum.
\end{enumerate}
Because all events are computed from the same pre-feedback state and summed
linearly, the update is independent of deposition order and handles overlapping
SN kernels deterministically.

\section{Alternative Hopkins Variants}
This section summarizes other variants described in Hopkins (2025) and earlier
FIRE implementations, which can be substituted into the coupling step above.

\subsection{Terminal-Momentum Variants}
The baseline uses
\begin{equation}
p_{0} = \min\left[p_{\rm egy},\,p_{\rm term}\right],
\end{equation}
with $p_{\rm term}$ taken from a density/metallicity scaling and $p_{\rm egy}$
from energy conservation. The paper discusses three practical choices for the
velocity dependence in $p_{\rm term}$:
\begin{enumerate}
\item \textbf{Fixed $F_{v}=1$:} ignore velocity dependence so
$p_{\rm term} = p_{t,0}\,F_{n}(\langle n\rangle)\,F_{Z}(\langle Z\rangle)$.
This matches most legacy sub-grid models and is nearly identical to the more
complex forms when $\langle v_{r}\rangle \ge 0$.
\item \textbf{Total-momentum change:} define $F_{v}$ using the analytic
terminal-momentum scaling in a background flow (Appendix~B.1 of Hopkins 2025),
which suppresses $p_{\rm term}$ for strong outflows and can increase it in
converging flows.
\item \textbf{$\Delta$-momentum change:} define $F_{v}$ using the difference
between the post-SN state and the no-SN evolution (Appendix~B.2). This tends to
avoid very large boosts in strongly converging flows and yields behavior closer
to $F_{v}=1$ for typical conditions.
\end{enumerate}
In all cases the coupling remains conservative; the choice primarily changes
how the unresolved work scales with $\langle v_{r}\rangle$.

\subsection{Energy vs. Momentum Deposition Limits}
The algorithm can be specialized to two limiting regimes by fixing $\chi$:
\begin{enumerate}
\item \textbf{Energy-conserving limit:} set $\chi=1$, so
$p_{0}=p_{\rm egy}$ and the coupled momentum is determined entirely by
energy conservation.
\item \textbf{Terminal-momentum limit:} fix $p_{0}=p_{\rm term}$ (equivalently
$\chi=p_{\rm term}/p_{\rm egy}$), representing unresolved cooling losses.
\end{enumerate}

\subsection{Directional vs. Isotropic Coupling}
Earlier FIRE implementations applied the solution ``in cones'' by computing
the energy-conserving factor $\psi$ (and hence $p_{0}$) separately for each
neighbor. This better preserves local anisotropy but breaks exact isotropy and
can up-weight the tails of the kernel. The Hopkins (2025) formulation uses a
single $\psi$ and $\chi$ per event with vector weights $\bar{\mathbf w}_{b}$,
ensuring exact isotropy and conservation. On a Cartesian mesh, the isotropic
version corresponds to the weight construction in Section~3, while the cone
version would apply the same weights but recompute $\psi$ using each
neighbor's local state.

\section{Energy Conservation and Galilean Invariance}
The Hopkins coupling enforces an exact energy conservation equation written in
terms of star-frame relative velocities. For a single event $a$ with neighbor
cells $b$, the discrete kinetic-energy balance can be written as (Hopkins 2025,
Appendix~A):
\begin{equation}
\epsilon_{a} = \sum_{b}\left[\frac{|\Delta \mathbf p_{b}|^{2}}{2 m_{b}(1+\mu_{b})}
+ \frac{\mathbf v_{ba}\cdot\Delta \mathbf p_{b}}{1+\mu_{b}}\right],
\end{equation}
with $\mu_{b}\equiv \Delta m_{b}/m_{b}$ and $\mathbf v_{ba}=\mathbf v_{b}-\mathbf v_{a}$.
Because $\mathbf v_{ba}$ is invariant under a uniform boost, this equation is
Galilean invariant. The coupled momentum $p_{0}$ is then chosen (possibly capped
by $p_{\rm term}$) so that the energy partition between kinetic and non-kinetic
reservoirs satisfies the same invariant relation. Any velocity dependence in
$p_{\rm term}$ should be expressed using star-frame or kernel-averaged relative
velocities to preserve invariance.

In contrast, the Quokka SN limiter operates on mesh-frame conserved variables
after summing all events in a cell, solving for a scaling $\lambda$ so that
$e_{\rm int}+|\mathbf p|^{2}/(2\rho)=e$ is satisfied in the lab frame. That
constraint is not written in terms of star-frame relative velocities, so the
resulting $\lambda$ can depend on the chosen inertial frame when multiple SNe
contribute to the same cell. A Galilean-invariant limiter would need
to enforce the constraint in a comoving frame (e.g., a kernel-averaged gas
frame) and then transform back, or avoid a post-hoc limiter by using the
Hopkins-style energy-conserving coupling per event.

\section{Implementation Notes}
\begin{itemize}
\item The weight construction uses only relative positions, so it is translation
invariant. The dipole removal ensures exact momentum conservation.
\item For AMR, use the local cell sizes in $A_{\rm proj}$ and the local $V_{b}$ in
all mass and energy updates.
\item The kernel radius should be at most the ghost-zone depth used in the
hydrodynamic solver to avoid extra communication steps.
\item The total coupled energy per event is not a fixed constant unless
$\mathbf v_{b}=\mathbf v_{a}$; the available energy is
$\mathcal E_{a}=E_{\rm ej}+\tfrac{1}{2}m_{\rm ej}\sum_{b}w^{\prime}_{b}|\mathbf v_{b}-\mathbf v_{a}|^{2}$.
In the terminal-momentum limit, part of $\mathcal E_{a}$ is assigned to the
non-kinetic reservoir $U_{a}$, so the explicit gas energy increment can be smaller.
\end{itemize}

\end{document}
